\documentclass[12pt,a4paper]{article}
\usepackage[utf8]{vietnam}
\usepackage{amsmath,amssymb,mathtools}
\usepackage{bm}
\usepackage{geometry}
\geometry{margin=2cm}
\usepackage{booktabs}
\usepackage{graphicx}
\usepackage{setspace}
\usepackage{hyperref}
\usepackage{parskip}
\usepackage{natbib}
\usepackage{float}
\usepackage{caption}
\usepackage{subcaption}

\onehalfspacing

% Custom math commands for readability
\newcommand{\E}{\mathbb{E}}
\newcommand{\V}{\mathbb{V}}
\newcommand{\R}{\mathbb{R}}
\newcommand{\N}{\mathcal{N}}
\newcommand{\CDF}{\hat{F}}
\newcommand{\CDFinv}{\hat{F}^{\,-1}}
\newcommand{\diff}{\mathop{}\!\Delta}
\newcommand{\abs}[1]{\lvert #1 \rvert}

\title{Phụ lục: Phương pháp sinh dữ liệu tổng hợp và đánh giá chất lượng}
\author{}
\date{}

\begin{document}
\maketitle

\section{Phương pháp sinh dữ liệu tổng hợp}

\subsection{Bối cảnh và thách thức}

Nghiên cứu khảo sát hành vi người tiêu dùng thế hệ Z trên TikTok Shop 
thu thập được 66 phản hồi, trong đó chỉ 39 quan sát có dữ liệu đầy đủ 
trên toàn bộ 49 biến Likert (thang đo 1--5). 
Với quy mô mẫu nhỏ như vậy, các phân tích thống kê đa biến 
như mô hình cấu trúc tuyến tính (SEM) hay phân tích nhân tố khẳng định (CFA) 
sẽ có độ mạnh thống kê (statistical power) thấp, 
dẫn đến nguy cơ cao mắc sai lầm loại II 
(không bác bỏ giả thuyết khống khi nó sai). 

Để khắc phục hạn chế này, chúng tôi áp dụng phương pháp sinh dữ liệu tổng hợp 
(synthetic data generation) nhằm mở rộng tập dữ liệu lên 466 quan sát 
(66 gốc + 400 tổng hợp), phục vụ cho các phân tích thăm dò 
(exploratory analysis) và kiểm tra độ vững của mô hình (robustness check).

\subsection{Lựa chọn phương pháp: Gaussian Copula kết hợp Conditional Bootstrap}

Chúng tôi lựa chọn phương pháp \textbf{Gaussian Copula} kết hợp 
\textbf{Conditional Bootstrap} dựa trên bốn lý do sau:

\subsubsection{(1) Bảo toàn cấu trúc tương quan giữa các biến}

Trong nghiên cứu hành vi, các khái niệm tiềm ẩn (constructs) như 
``Mua bốc đồng'', ``Hối hận sau mua'', ``Ý định trả hàng'' 
có mối tương quan chặt chẽ với nhau. 
Nếu sinh dữ liệu từng cột riêng rẽ, cấu trúc tương quan này sẽ bị phá vỡ, 
khiến dữ liệu tổng hợp mất đi ý nghĩa nghiên cứu.

Gaussian Copula giải quyết vấn đề này bằng cách:
\begin{enumerate}
    \item Ước lượng ma trận tương quan $\bm{\Sigma}$ từ dữ liệu gốc
    \item Sinh các mẫu từ phân phối chuẩn đa biến $\N(\bm{0}, \bm{\Sigma})$
    \item Chiếu về phân phối gốc thông qua hàm phân phối tích lũy nghịch đảo
\end{enumerate}

Cụ thể, gọi $\CDF_j$ là hàm phân phối tích lũy (CDF) thực nghiệm 
của biến $X_j$ và $\Phi$ là CDF của phân phối chuẩn hóa. 
Quy trình chuyển đổi như sau:

\begin{align}
    U_{ij} &= \CDF_j(X_{ij})                    && \text{(chuyển về phân phối đều)}   \label{eq:uniform} \\
    Z_{ij} &= \Phi^{-1}(U_{ij})                  && \text{(chuyển về phân phối chuẩn)} \label{eq:normal} \\
    \bm{\Sigma} &= \operatorname{Corr}(\bm{Z})   && \text{(ma trận tương quan)}         \label{eq:corr} \\
    \bm{Z}_{\text{new}} &\sim \N(\bm{0}, \bm{\Sigma}) && \text{(sinh mẫu mới)}          \label{eq:sampling} \\
    \bm{U}_{\text{new}} &= \Phi(\bm{Z}_{\text{new}})  && \text{(quay về phân phối đều)} \label{eq:back} \\
    X_{ij}^{\text{(new)}} &= \CDFinv_j\bigl(U_{ij}^{\text{(new)}}\bigr) && \text{(chiếu về phân phối gốc)} \label{eq:inverse}
\end{align}

Vì dữ liệu Likert là rời rạc (discrete), chúng tôi sử dụng kỹ thuật 
``jittering'' — thêm nhiễu ngẫu nhiên vào các giá trị trùng 
trong quá trình ước lượng CDF — để tránh hiện tượng 
các giá trị bằng nhau gây ra ma trận tương quan suy biến.

\subsubsection{(2) Bảo toàn phân phối xác suất của từng biến}

Khác với các phương pháp tham số (parametric) giả định phân phối chuẩn, 
Gaussian Copula sử dụng CDF thực nghiệm (empirical CDF), 
cho phép tái tạo chính xác mọi dạng phân phối — 
kể cả phân phối lệch (skewed) hay đa mode. 
Điều này đặc biệt quan trọng vì dữ liệu Likert trong nghiên cứu 
thường có phân phối lệch trái hoặc lệch phải.

\subsubsection{(3) Tôn trọng cấu trúc logic của bảng hỏi}

Khảo sát có cấu trúc nhánh (skip logic): 
người chưa từng mua hàng qua livestream (``Chưa'') 
sẽ không trả lời các câu hỏi về cảm nhận sau mua, 
dẫn đến dữ liệu khuyết (missing) có hệ thống.

Chúng tôi kết hợp Conditional Bootstrap để xử lý:
\begin{itemize}
    \item Tính tỷ lệ ``Có'' / ``Chưa'' từ 66 mẫu gốc ($78\%$ -- $22\%$)
    \item Sinh 310 dòng ``Có'' từ Gaussian Copula (đầy đủ 49 biến Likert)
    \item Sinh 90 dòng ``Chưa'' từ Bootstrap thuần túy (\texttt{NaN} ở 49 biến Likert)
\end{itemize}

Phương pháp này đảm bảo dữ liệu tổng hợp không vi phạm logic 
nội tại của bảng hỏi và giữ nguyên tỷ lệ tổng thể.

\subsubsection{(4) Khả năng mở rộng với cỡ mẫu nhỏ}

Với chỉ 39 quan sát hoàn chỉnh, các phương pháp sinh dữ liệu 
phức tạp hơn như mô hình sinh sâu (deep generative models — 
GAN, VAE, diffusion models) không khả thi do thiếu dữ liệu huấn luyện. 
Gaussian Copula chỉ yêu cầu ước lượng:
\begin{itemize}
    \item Các phân phối biên (marginal distributions) — 39 điểm đủ cho CDF thực nghiệm
    \item Ma trận tương quan — với 49 biến, cần ít nhất 
    $\frac{49 \times 50}{2} = 1225$ tham số, nhưng với 39 mẫu, 
    ta bị underdetermined (ít mẫu hơn số tham số)
\end{itemize}

Hạn chế về cỡ mẫu được giải quyết bằng cách regularize ma trận tương quan 
(phương trình~\ref{eq:regularize}):

\begin{equation}
    \bm{\Sigma}_{\text{reg}} = \bm{\Sigma} + \lambda \bm{I}, 
    \qquad \lambda = \max\!\bigl(-\lambda_{\min}(\bm{\Sigma}),\,0\bigr) + \varepsilon
    \label{eq:regularize}
\end{equation}

trong đó $\varepsilon = 10^{-6}$ là hằng số điều chỉnh nhằm đảm bảo 
ma trận $\bm{\Sigma}_{\text{reg}}$ xác định dương (positive definite).

\subsection{Tổng quan quy trình}

Bảng~\ref{tab:pipeline} tóm tắt toàn bộ quy trình xử lý:

\begin{table}[H]
\centering
\caption{Quy trình sinh dữ liệu tổng hợp}
\label{tab:pipeline}
\begin{tabular}{cll}
\toprule
Bước & Xử lý & Chi tiết \\
\midrule
1 & Tiền xử lý & Làm sạch, rename cột, loại bỏ chuỗi rỗng \\
2 & Tách nhóm & 39 dòng đầy đủ (nhóm Copula) + 27 dòng khuyết \\
3 & Ước lượng & Phân phối biên + Ma trận tương quan từ 39 dòng \\
4 & Copula & Sinh 400 mẫu MVN, chiếu về Likert 1--5 \\
5 & Bootstrap & Sinh categorical theo tỷ lệ gốc (310 Có / 90 Chưa) \\
6 & Kiểm định & KS test, so sánh tương quan, phân phối tần số \\
\bottomrule
\end{tabular}
\end{table}

\newpage
\section{Đánh giá chất lượng bộ dữ liệu tổng hợp}

\subsection{Kết quả kiểm định}

Chúng tôi thực hiện bốn nhóm kiểm định để đánh giá chất lượng 
của 400 dòng dữ liệu tổng hợp so với 39 dòng gốc hoàn chỉnh:

\subsubsection{Kiểm định Kolmogorov--Smirnov (KS test)}

Kết quả KS test 2-mẫu trên từng cột Likert cho thấy:

\begin{table}[H]
\centering
\caption{Kết quả KS test trên 49 biến Likert}
\label{tab:kstest}
\begin{tabular}{lcc}
\toprule
Chỉ số & Giá trị & Đánh giá \\
\midrule
Số cột PASS ($p > 0{,}05$) & 49/49 & Xuất sắc \\
Mean KS statistic & 0,034 & Rất thấp \\
Min KS statistic & 0,009 & --- \\
Max KS statistic & 0,089 & --- \\
Mean $p$-value & 0,998 & --- \\
Min $p$-value & 0,924 & --- \\
\bottomrule
\end{tabular}
\end{table}

Với $100\%$ số cột vượt qua KS test ở mức ý nghĩa $5\%$, 
phân phối của các biến Likert tổng hợp là đồng nhất về mặt thống kê 
với phân phối gốc. 

\subsubsection{Sai lệch tham số thống kê}

\begin{table}[H]
\centering
\caption{Sai lệch tham số giữa dữ liệu gốc và tổng hợp}
\label{tab:param}
\begin{tabular}{lccc}
\toprule
Tham số & $\operatorname{Mean}\abs{\diff}$ & Giá trị lớn nhất & Đánh giá \\
\midrule
Trung bình (Mean) & 0,061 & 0,181 (thang 1--5) & Tốt \\
Độ lệch chuẩn (Std) & 0,036 & 0,106 & Rất tốt \\
Trung vị (Median) & 0,010 & 0,500 & Rất tốt \\
\bottomrule
\end{tabular}
\end{table}

Sai lệch trung bình $0{,}061$ trên thang đo 5 điểm tương đương 
$1{,}5\%$ toàn thang đo. 
Đây là mức sai lệch chấp nhận được với cỡ mẫu gốc nhỏ ($n = 39$).

\subsubsection{Bảo toàn cấu trúc tương quan}

Bảng~\ref{tab:corr} so sánh ma trận tương quan của 49 biến Likert:

\begin{table}[H]
\centering
\caption{So sánh ma trận tương quan}
\label{tab:corr}
\begin{tabular}{lcc}
\toprule
Chỉ số & Giá trị & Đánh giá \\
\midrule
$\operatorname{Mean}\abs{\diff}$ (toàn bộ) & 0,136 & Chấp nhận được \\
$\operatorname{Mean}\abs{\diff}$ (off-diagonal) & 0,139 & Chấp nhận được \\
Độ lệch chuẩn của $\diff$ (off-diagonal) & 0,075 & --- \\
$\max\abs{\diff}$ & 0,383 & --- \\
$\mathrm{Q}_{25}$ của $\diff$ & 0,084 & --- \\
$\mathrm{Q}_{75}$ của $\diff$ & 0,193 & --- \\
\bottomrule
\end{tabular}
\end{table}

\subsection{Phân tích các điểm yếu và hệ quả}

Mặc dù kết quả KS test và sai lệch tham số ở mức tốt, 
chúng tôi ghi nhận ba điểm yếu cần phân tích:

\subsubsection{Điểm yếu 1: Suy giảm tương quan nội bộ nhóm}

Ma trận tương quan trung bình trong từng nhóm biến 
(in-group correlation) có sự suy giảm đáng kể.

\begin{table}[H]
\centering
\caption{So sánh tương quan trung bình trong từng nhóm}
\label{tab:groupcorr}
\begin{tabular}{lccc}
\toprule
Nhóm & Gốc ($r$) & Tổng hợp ($r$) & Chênh lệch \\
\midrule
Chất lượng thông tin & 0,715 & 0,505 & $-0{,}210$ \\
Tương tác Streamer & 0,604 & 0,393 & $-0{,}211$ \\
KOL/KOC & 0,760 & 0,584 & $-0{,}176$ \\
Dòng chảy/Cuốn hút & 0,766 & 0,627 & $-0{,}138$ \\
Mua bốc đồng & 0,705 & 0,571 & $-0{,}135$ \\
Hối hận sau mua & 0,650 & 0,477 & $-0{,}173$ \\
Ý định trả hàng & 0,754 & 0,531 & $-0{,}223$ \\
Truyền miệng tiêu cực & 0,625 & 0,496 & $-0{,}129$ \\
\bottomrule
\end{tabular}
\end{table}

\textbf{Nguyên nhân:} 
Hiện tượng suy giảm tương quan nội bộ nhóm (gọi là ``correlation shrinkage'') 
bắt nguồn từ việc ma trận tương quan được ước lượng trên 
chỉ 39 quan sát. Với $p = 49$ biến, số tham số cần ước lượng là 
$p(p-1)/2 = 1176$, trong khi chỉ có 39 mẫu, dẫn đến 
tình trạng ``underdetermined'' — số tham số vượt xa số quan sát.
Ma trận tương quan ước lượng được có phương sai lớn 
và một số trị riêng bằng $0$ (singular). 
Mặc dù đã regularize (phương trình~\ref{eq:regularize}), 
quá trình sinh mẫu từ copula không thể tái tạo hoàn hảo 
các tương quan cao do thiếu thông tin.

\textbf{Hệ quả đối với nghiên cứu:}

\begin{enumerate}
    \item \textbf{Hệ số Cronbach's Alpha bị ước lượng thấp hơn thực tế:}
    Vì $\alpha$ phụ thuộc trực tiếp vào tương quan trung bình giữa các biến 
    trong cùng một thang đo, dữ liệu tổng hợp sẽ cho giá trị $\alpha$ 
    thấp hơn so với thực tế. Điều này có thể dẫn đến kết luận sai lầm rằng 
    thang đo không đạt độ tin cậy (ngưỡng thường là $\alpha > 0{,}7$), 
    trong khi thực tế thang đo hoàn toàn tốt.
    
    \item \textbf{Phương sai trích trung bình (AVE) bị suy giảm:}
    Trong phân tích CFA và PLS-SEM, AVE được tính từ các trọng số nhân tố 
    (factor loadings). Khi tương quan nội bộ suy giảm, 
    các trọng số nhân tố cũng giảm theo, kéo AVE xuống dưới ngưỡng $0{,}5$.
    Điều này có thể khiến nhà nghiên cứu kết luận 
    ``giá trị hội tụ (convergent validity) không đạt'' 
    trong khi thực tế giá trị hội tụ là đạt yêu cầu.
    
    \item \textbf{Hệ số hồi quy giữa các nhân tố bị thiên lệch:}
    Trong mô hình SEM, tương quan giữa các nhân tố 
    (inter-construct correlation) được bảo toàn tốt hơn tương quan nội bộ, 
    nhưng khi AVE bị suy giảm, các hệ số đường dẫn (path coefficients) 
    trong mô hình cấu trúc có thể bị ước lượng với sai số lớn hơn.
    Các kiểm định giả thuyết (hypothesis testing) dựa trên 
    bootstrap confidence interval trong PLS-SEM vẫn có giá trị 
    tham khảo nhưng cần thận trọng khi diễn giải.
\end{enumerate}

\textbf{Khuyến nghị giảm thiểu:}
\begin{itemize}
    \item Sử dụng dữ liệu tổng hợp chủ yếu cho phân tích thăm dò 
    (exploratory factor analysis — EFA) và kiểm tra độ vững.
    \item Khi đánh giá độ tin cậy thang đo trên dữ liệu tổng hợp, 
    cần điều chỉnh kỳ vọng: chấp nhận ngưỡng $\alpha < 0{,}7$ 
    nếu dữ liệu gốc cho thấy thang đo tốt.
    \item \textbf{KHÔNG} sử dụng dữ liệu tổng hợp để thay thế dữ liệu gốc 
    trong các kiểm định giả thuyết chính thức (confirmatory analysis). 
    Dữ liệu tổng hợp chỉ nên dùng để tăng cường (augment) 
    chứ không thay thế dữ liệu thật.
\end{itemize}

\subsubsection{Điểm yếu 2: Sai lệch có hệ thống ở một số biến}

Một số biến có sai lệch trung bình lớn hơn $0{,}15$ 
(\texttt{Regret\_Purchase\_Decision}: $0{,}18$; 
\texttt{Frequent\_Unplanned\_Purchases}: $0{,}17$; 
\texttt{See\_Buy\_Behavior}: $0{,}17$). 
Các biến này đều thuộc nhóm ``Mua bốc đồng'' và ``Hối hận sau mua''.

\textbf{Nguyên nhân:} 
Các biến này có phân phối lệch (skewed) mạnh và 
phương sai lớn trong dữ liệu gốc ($\sigma > 1{,}3$), 
khiến CDF thực nghiệm kém ổn định hơn so với các biến 
có phương sai nhỏ. Với 39 mẫu, 
các giá trị ngoại lệ (outliers) có ảnh hưởng lớn đến 
hình dạng phân phối.

\textbf{Hệ quả:}
\begin{itemize}
    \item Các phân tích mô tả (descriptive statistics) trên dữ liệu tổng hợp 
    sẽ có sai số nhỏ nhưng có hệ thống ở các biến này.
    \item Trong mô hình hồi quy, nếu các biến này là biến phụ thuộc,
    ước lượng hệ số có thể bị chệch.
    \item Tuy nhiên, với $\abs{\diff_{\text{mean}}} = 0{,}18$ trên thang 5 điểm, 
    sai lệch này vẫn nằm trong khoảng $\pm 0{,}5$ điểm — 
    không làm thay đổi ý nghĩa thực tiễn của kết quả 
    (ví dụ: mean gốc $3{,}0$ ``trung bình'' vẫn là mean tổng hợp $2{,}82$ ``trung bình'').
\end{itemize}

\subsubsection{Điểm yếu 3: Không bắt được cấu trúc đa mode hiếm}

Với 39 mẫu, CDF thực nghiệm có độ phân giải thấp. 
Nếu một giá trị Likert chỉ xuất hiện 1--2 lần trong 39 mẫu, 
thì CDF thực nghiệm ước lượng xác suất của nó với sai số lớn. 
Copula sinh ra phân phối trơn (smooth) hơn, 
làm mất các mode hiếm.

Ví dụ: nếu giá trị ``5'' xuất hiện $2/39 \approx 5\%$ trong dữ liệu gốc, 
Copula có thể sinh ra nó với tần suất $3$--$8\%$ 
(sai số tương đối đến $60\%$).

\textbf{Hệ quả:}
\begin{itemize}
    \item Các biến có ít giá trị ngoại cùng (ceiling/floor effects) 
    sẽ bị tái tạo kém chính xác nhất.
    \item Khi dùng dữ liệu tổng hợp để nghiên cứu 
    hành vi cực đoan (ví dụ: người rất hối hận với score 5 ở mọi biến), 
    cần kiểm tra lại bằng dữ liệu gốc.
    \item Đây là hạn chế cố hữu của mọi phương pháp tham số hóa 
    với cỡ mẫu nhỏ, không riêng gì Copula.
\end{itemize}

\subsection{Kết luận về chất lượng}

\begin{table}[H]
\centering
\caption{Đánh giá tổng thể chất lượng dữ liệu tổng hợp}
\label{tab:summary}
\begin{tabular}{lp{6cm}p{4cm}}
\toprule
Tiêu chí & Kết quả & Mức độ \\
\midrule
Phân phối đơn biến & KS test: 49/49 PASS ($p > 0{,}05$) & Xuất sắc \\
Sai lệch trung bình & $\operatorname{Mean}\abs{\diff} = 0{,}061$ ($1{,}5\%$ thang đo) & Tốt \\
Sai lệch độ lệch chuẩn & $\operatorname{Mean}\abs{\diff} = 0{,}036$ & Rất tốt \\
Tương quan giữa nhóm & $\operatorname{Mean}\abs{\diff} = 0{,}10$--$0{,}15$ & Khá \\
Tương quan trong nhóm & Giảm $13$--$22\%$ so với gốc & Yếu \\
Skip logic & $100\%$ chính xác & Xuất sắc \\
\bottomrule
\end{tabular}
\end{table}

Dữ liệu tổng hợp đạt chất lượng \textbf{chấp nhận được} 
cho các phân tích thăm dò (EFA, độ vững mô hình) 
nhưng \textbf{không thay thế} dữ liệu gốc trong các kiểm định 
giả thuyết chính thức. 
Điểm yếu chính là suy giảm tương quan nội bộ nhóm 
do cỡ mẫu gốc quá nhỏ ($n = 39$), 
dẫn đến nguy cơ đánh giá thấp độ tin cậy thang đo. 
Khuyến nghị sử dụng dữ liệu tổng hợp như một công cụ hỗ trợ 
(tăng cường, kiểm tra độ vững), không phải nguồn dữ liệu chính thức.

\end{document}
